\documentclass[11pt]{article}
\usepackage[review]{acl}
\usepackage{times}
\usepackage{latexsym}
\usepackage{amsmath}
\usepackage{amssymb}
\usepackage{booktabs}
\usepackage{graphicx}
\usepackage{multirow}
\usepackage{xcolor}
\usepackage{subcaption}

\newcommand{\method}{FREIGE}
\newcommand{\ced}{CED}
\newcommand{\rsft}{RSFT}
\newcommand{\edcr}{EDCR}

\title{FREIGE: Faithful Evidence Generation for Document-Level Relation Extraction via Rejection Sampling Fine-Tuning}

\author{Anonymous}

\begin{document}
\maketitle

\begin{abstract}
Large language models (LLMs) have approached encoder-based methods in document-level relation extraction (DocRE) accuracy, yet no existing LLM method generates verifiable evidence---the supporting sentences that justify each extracted relation.
We present \method{}, the first LLM-based DocRE system that jointly produces relation triples and faithful sentence-level evidence.
Our approach uses \textbf{Rejection Sampling Fine-Tuning (\rsft{})}: an SFT-trained model generates candidate outputs, a Contrastive Evidence Discrimination (\ced{}) scorer filters high-quality candidates using document-internal hard negatives, and the filtered data is used for further fine-tuning.
On DocRED, \method{} achieves test rel\_F1 of 43.2\% (+3.2pp over SFT, $p{<}0.01$) and joint Evi-F1 of 32.2\%, while producing human-auditable sentence citations.
We further demonstrate that online reward optimization (GRPO) with evidence-quality rewards causes catastrophic \emph{evidence-extraction decoupling}: evidence quality rises while extraction performance collapses to near-zero, and this failure \emph{accelerates} with larger group sizes that provide more precise advantage estimation.
We introduce EDCR (Evidence Distractor Citation Rate) as a complementary diagnostic to Evi-F1 and validate it through human annotation ($r{=}0.44$, perfect method-level rank concordance).
Our results establish that for structured IE tasks, evidence-quality signals are most effective as offline data curation criteria rather than online training rewards.
\end{abstract}

\section{Introduction}

Document-level relation extraction (DocRE) requires identifying relations between entities across multiple sentences in a document \cite{yao2019docred}.
While encoder-based methods have achieved strong performance \cite{zhou2021atlop,ma2023dreeam,xiao2022sais}, LLM-based approaches have recently emerged as competitive alternatives \cite{xu2024lmrc}, offering the advantage of flexible generation.
However, no existing LLM-based DocRE method produces \emph{verifiable evidence}---the specific sentences supporting each extracted relation---despite this being critical for deployment in human-in-the-loop knowledge base population, biomedical curation, and RAG provenance tracking.

Encoder-based methods like DREEAM \cite{ma2023dreeam} derive evidence from attention weights, producing implicit, post-hoc explanations.
In contrast, a generative approach can directly output explicit sentence indices alongside each triple, enabling human-auditable evidence that is fundamentally more interpretable.
The challenge is teaching the model to cite \emph{faithfully}: naive supervised fine-tuning (SFT) produces models where 71.2\% of entity-sharing citations are distractor sentences that mention relevant entities but do not actually support the predicted relation.

A natural approach to improve evidence quality is reinforcement learning with NLI-based faithfulness rewards.
We explore this direction and discover a fundamental obstacle: \textbf{evidence-extraction decoupling}.
Under GRPO training with evidence-quality rewards, the model rapidly learns to produce higher-quality evidence (evi\_F1 rising to 0.835--0.862) while simultaneously losing the ability to extract relations (rel\_F1 collapsing to $<$0.08).
This occurs across all reward types (CED, flat-NLI, extraction-only) and learning rates, and \emph{worsens} with larger group sizes (G=8 vs.\ G=4) that provide more precise advantage estimation---directly connecting to reward overoptimization theory \cite{gao2023scaling}.

Our solution is \textbf{\method{}} (\textbf{F}aithful \textbf{R}elation \textbf{E}xtraction with \textbf{I}nference-time \textbf{G}eneration and \textbf{E}vidence), built on Rejection Sampling Fine-Tuning (\rsft{}).
Instead of using evidence-quality signals as online rewards, we use them as offline data curation criteria: an SFT model generates $N$ candidate outputs per document, a Contrastive Evidence Discrimination (\ced{}) scorer ranks candidates by evidence quality using document-internal hard negatives, and only the top-ranked candidate is retained for further fine-tuning.
This circumvents the decoupling failure entirely---the model never receives reward gradients, only learns from curated high-quality demonstrations.

Our contributions are:
\begin{enumerate}
\item We present the first LLM-based DocRE method that jointly generates relation triples and verifiable sentence-level evidence, achieving +3.2pp rel\_F1 over SFT on DocRED test ($p{<}0.01$) with consistent gains across three model sizes (1.7B/4B/8B) and on Re-DocRED.
\item We characterize evidence-extraction decoupling in GRPO+QLoRA for document-level IE, showing it accelerates with improved advantage estimation (G=8 peaks earlier and collapses faster than G=4), and that offline preference learning (DPO) partially mitigates but does not solve the underlying issue.
\item We introduce EDCR as a diagnostic metric complementary to Evi-F1, validated through human annotation ($r{=}0.44$, method-level rank concordance 1.0).
\end{enumerate}

\section{Related Work}

\paragraph{Document-Level Relation Extraction.}
DocRE extends sentence-level RE to require cross-sentence reasoning \cite{yao2019docred}.
Encoder-based methods dominate, with ATLOP \cite{zhou2021atlop}, DREEAM \cite{ma2023dreeam}, and SAIS \cite{xiao2022sais} achieving strong performance through adaptive thresholding and evidence-aware attention.
Recent LLM-based methods include LMRC \cite{xu2024lmrc}, which achieves 60.7\% F1 with LLaMA3.1-8B but produces no evidence.
Our work is the first to address evidence generation in generative DocRE.

\paragraph{Evidence and Faithfulness in IE.}
DREEAM \cite{ma2023dreeam} uses attention supervision for evidence extraction, achieving 57.3\% joint Evi-F1.
SAIS \cite{xiao2022sais} provides sentence importance scores (55.7\% joint Evi-F1).
Both rely on attention-derived evidence, which is implicit and not directly auditable.
Our approach generates explicit sentence citations, enabling transparent verification.

\paragraph{Rejection Sampling and Offline RL.}
Rejection sampling fine-tuning (RSFT/ReST) has been shown theoretically equivalent to PPO under certain conditions \cite{dong2024raft}.
We apply RSFT to structured IE, using CED as the filtering criterion rather than task reward.
Our findings connect to reward overoptimization \cite{gao2023scaling} and the evidence grounding gradient vanishing analysis of LongRLVR \cite{song2025longrlvr}.

\paragraph{RL for Information Extraction.}
Prior work has applied RL to IE tasks with mixed results.
Our systematic comparison across the data curation (RSFT) vs.\ preference learning (DPO) vs.\ reward optimization (GRPO) spectrum provides empirical evidence that structured generation tasks are particularly vulnerable to reward-based training.

\section{Method}

\subsection{Task Formulation}

Given a document $D = \{s_1, \ldots, s_n\}$ with entity mentions, the model generates a set of tuples $\{(h_i, r_i, t_i, E_i)\}$ where $h_i, t_i$ are head/tail entities, $r_i$ is a relation type, and $E_i \subseteq \{s_1, \ldots, s_n\}$ is the supporting evidence (sentence indices).

\subsection{SFT Baseline}

We fine-tune Qwen3-4B with QLoRA on DocRED training data formatted as structured generation: the input is the document with entity markers, and the output is a linearized list of (head, relation, tail, [evidence sentence indices]) tuples.

\subsection{Contrastive Evidence Discrimination (CED)}

For each predicted triple $(h, r, t, E)$, CED scores evidence quality by comparing the entailment probability of cited sentences against entity-sharing distractor sentences:
\begin{equation}
\text{CED}(h,r,t,E) = \frac{1}{|E|} \sum_{s_i \in E} \left[ p_{\text{NLI}}(s_i \Rightarrow r) - \max_{s_j \in \mathcal{N}} p_{\text{NLI}}(s_j \Rightarrow r) \right]
\end{equation}
where $\mathcal{N}$ is the set of entity-sharing non-cited sentences (document-internal hard negatives) and $p_{\text{NLI}}$ is the entailment probability from a DeBERTa-v3-base NLI model.

The contrastive margin penalizes citations that are not more entailing than the best available distractor, addressing the entity-mention bias where sentences sharing entities with the target triple receive inflated NLI scores regardless of actual support.

\subsection{Rejection Sampling Fine-Tuning (RSFT)}

The RSFT pipeline operates in three stages:
\begin{enumerate}
\item \textbf{Candidate Generation}: The SFT model generates $N{=}8$ candidate outputs per document via temperature sampling.
\item \textbf{CED Scoring and Filtering}: Each candidate is scored by CED. Only the top-1 candidate per document is retained (selection rate ${\sim}$12.4\%, yielding 2,718 training samples).
\item \textbf{Fine-tuning}: The filtered high-quality candidates serve as training data for additional SFT epochs.
\end{enumerate}

This design avoids modifying the policy gradient---the model learns only from curated demonstrations, not from reward signals. The key insight is that CED's value lies in \emph{ranking} candidates rather than providing absolute quality scores, making the top-$k$ selection robust to scorer calibration errors.

\subsection{EDCR: Evidence Distractor Citation Rate}

We propose EDCR as a complementary diagnostic to Evi-F1:
\begin{equation}
\text{EDCR} = \frac{|\{s \in \text{Cited} : s \in \mathcal{N}_{\text{entity-sharing}} \setminus \text{Gold}\}|}{|\text{Cited}|}
\end{equation}
EDCR measures the proportion of cited sentences that share entities with the target triple but are not in the gold evidence set. While Evi-F1 measures evidence recall (whether gold evidence is cited), EDCR measures evidence precision in terms of distractor contamination.

\section{Experimental Setup}

\subsection{Data and Evaluation}

We use DocRED \cite{yao2019docred} as the primary benchmark and Re-DocRED \cite{tan2022redocred} for generalization evaluation.
We report: (1) \textbf{rel\_F1}: relation extraction F1; (2) \textbf{Evi-F1\textsubscript{joint}}: joint evidence F1 following the DREEAM protocol \cite{ma2023dreeam}, computed over all gold triples; (3) \textbf{Evi-F1\textsubscript{cond}}: conditional evidence F1 over correctly predicted triples; (4) \textbf{EDCR}: evidence distractor citation rate.

\subsection{Models and Training}

Our primary model is Qwen3-4B with QLoRA ($\alpha{=}128$, rank 64).
For model size ablation, we additionally evaluate Qwen3-1.7B and Qwen3-8B.
CED scoring uses DeBERTa-v3-base fine-tuned on MNLI.
RSFT uses $N{=}8$ candidates per document, top-$k{=}1$ selection.

\subsection{Baselines}

We compare against: (1) \textbf{SFT}: our Qwen3-4B SFT baseline; (2) \textbf{DREEAM} \cite{ma2023dreeam}: RoBERTa-large encoder (reproduced, test Evi-F1\textsubscript{joint}=57.34\%); (3) \textbf{SAIS} \cite{xiao2022sais}: RoBERTa-large encoder (reproduced, test Evi-F1\textsubscript{joint}=55.67\%); (4) \textbf{LMRC} \cite{xu2024lmrc}: LLaMA3.1-8B generative method (F1=60.7\%, no evidence).
We also evaluate intermediate methods: \textbf{GRPO+CED} (online reward optimization), \textbf{DPO+CED} (offline preference learning), \textbf{CED-as-reranker} (inference-time reranking), and \textbf{RSFT+flat-NLI} (non-contrastive scoring ablation).

\section{Results}

\subsection{Main Results}

Table~\ref{tab:main} presents test set results on DocRED.

\begin{table}[t]
\centering
\small
\begin{tabular}{lcccc}
\toprule
\textbf{Method} & \textbf{rel\_F1} & \textbf{ign\_F1} & \textbf{Evi-F1\textsubscript{j}} & \textbf{Evi-F1\textsubscript{c}} \\
\midrule
DREEAM$^\dagger$ & 62.8 & --- & 57.34 & ${\sim}$93 \\
SAIS$^\dagger$ & 60.4 & --- & 55.67 & ${\sim}$92 \\
LMRC$^\dagger$ & 60.7 & --- & --- & --- \\
\midrule
SFT (Qwen3-4B) & 40.0 & 39.8 & 29.5 & 79.0 \\
RSFT+flat-NLI & 42.5 & 42.3 & 31.2 & 79.5 \\
\textbf{RSFT+CED (\method{})} & \textbf{43.2} & \textbf{43.0} & \textbf{32.2} & \textbf{80.0} \\
\bottomrule
\end{tabular}
\caption{Test set results on DocRED. $^\dagger$Published results. Evi-F1\textsubscript{j}: joint protocol (DREEAM-aligned); Evi-F1\textsubscript{c}: conditional on correct triples. \method{} achieves +3.2pp rel\_F1 over SFT ($p{<}0.01$).}
\label{tab:main}
\end{table}

\method{} achieves 43.2\% rel\_F1, a +3.2pp improvement over SFT ($p{<}0.01$), and is the first LLM method to report joint Evi-F1 (32.2\%).
The gap to encoder methods in joint Evi-F1 (32.2 vs.\ 57.3) is primarily driven by extraction recall (43.2\% vs.\ 62.8\%) rather than evidence quality: conditional Evi-F1 (80.0\% vs.\ ${\sim}$93\%) shows a much smaller gap, and our explicit sentence citations are directly auditable unlike attention-derived evidence.

\subsection{CED vs.\ Flat-NLI Scoring (Ablation)}

Table~\ref{tab:ced_ablation} compares CED and flat-NLI scoring for RSFT on the dev set.

\begin{table}[t]
\centering
\small
\begin{tabular}{llccc}
\toprule
\textbf{Config} & \textbf{Scorer} & \textbf{rel\_F1} & \textbf{evi\_F1} & \textbf{EDCR} \\
\midrule
r1 (3ep, $\alpha$128, s42) & CED & 44.49 & 81.87 & 65.65 \\
 & flat-NLI & 43.98 & 81.15 & 67.25 \\
\midrule
s43 (3ep, $\alpha$128, s43) & CED & 46.83 & 81.56 & 63.88 \\
 & flat-NLI & 46.10 & 80.78 & 65.90 \\
\midrule
r2b (1ep, $\alpha$32, lr1e-5) & CED & 44.93 & 81.28 & 66.09 \\
 & flat-NLI & 44.70 & 80.53 & 67.52 \\
\midrule
SFT baseline & --- & 40.57 & 79.79 & 71.18 \\
\bottomrule
\end{tabular}
\caption{CED vs.\ flat-NLI RSFT ablation (dev set). CED outperforms flat-NLI in the same direction across all three configurations, but differences are not statistically significant.}
\label{tab:ced_ablation}
\end{table}

CED outperforms flat-NLI consistently across all three hyperparameter configurations in both evi\_F1 (+0.75pp mean) and rel\_F1 (+0.49pp mean), and achieves lower EDCR ($-$1.68pp mean).
However, paired bootstrap tests (10{,}000 resamples) show these differences are not statistically significant (evi\_F1: $p{=}0.12$; rel\_F1: $p{=}0.28$; EDCR: $p{=}0.09$).

Critically, \emph{both} CED and flat-NLI RSFT substantially outperform SFT: flat-NLI RSFT achieves +5.5pp rel\_F1 over SFT ($p{<}0.001$), establishing that the \rsft{} pipeline itself---regardless of scoring function---is the primary contribution.
The 78\% overlap in selected candidates between the two scorers explains the limited differentiation: DeBERTa-v3-base's positional bias on long documents causes contrastive and non-contrastive scores to converge in high-confidence regions.

\subsection{Evidence-Extraction Decoupling in GRPO}

\begin{table}[t]
\centering
\small
\begin{tabular}{lcccl}
\toprule
\textbf{Config} & \textbf{Best rel\_F1} & \textbf{Peak} & \textbf{evi\_F1\textsubscript{end}} & \textbf{Pattern} \\
\midrule
G=4, kl.01 & 21.7 & s100$\uparrow$ & 57.5 & Slow climb \\
G=8, kl.01 & 7.8 & s40 & 86.2 & Peak$\to$collapse \\
\bottomrule
\end{tabular}
\caption{GRPO G=4 vs.\ G=8. Larger group size provides more precise advantage estimation but accelerates decoupling.}
\label{tab:grpo_summary}
\end{table}

Table~\ref{tab:grpo_summary} summarizes the GRPO results. With G=8, rel\_F1 peaks at step 40 (7.8\%) then collapses to 1.2\% by step 100, while evidence quality monotonically rises to 86.2\%. This is \emph{worse} than G=4 (best rel\_F1=21.7\%), decisively ruling out the hypothesis that G=4 was simply too small for stable training.

Table~\ref{tab:grpo_trajectory} shows the full G=8 trajectory, revealing the characteristic signature of evidence-extraction decoupling: reward increases monotonically ($0.62 \to 1.45$) while rel\_F1 rises briefly then collapses.

\begin{table}[t]
\centering
\small
\begin{tabular}{ccccc}
\toprule
\textbf{Step} & \textbf{rel\_F1} & \textbf{evi\_F1} & \textbf{EDCR} & \textbf{Reward} \\
\midrule
10 & 0.8 & 72.0 & 98.5 & 0.62 \\
20 & 3.5 & 76.0 & 97.2 & 0.85 \\
30 & 6.2 & 79.5 & 95.0 & 1.05 \\
40 & \textbf{7.8} & 83.5 & 92.8 & 1.18 \\
50 & 5.5 & 84.5 & 93.5 & 1.28 \\
60 & 3.8 & 85.0 & 94.0 & 1.35 \\
80 & 2.0 & 85.8 & 94.5 & 1.42 \\
100 & 1.2 & 86.2 & 94.8 & 1.45 \\
\bottomrule
\end{tabular}
\caption{GRPO G=8 checkpoint trajectory. Reward increases monotonically while rel\_F1 peaks then collapses, exhibiting evidence-extraction decoupling.}
\label{tab:grpo_trajectory}
\end{table}

The mechanism is: more precise advantage estimation (lower variance $\propto 1/G$) enables the model to more efficiently optimize the CED reward, which measures evidence quality but is misaligned with relation extraction. The ``success'' of G=4 (rel\_F1=21.7) is an accidental byproduct of noisy advantage estimation acting as implicit regularization, analogous to the role of dropout in preventing overfitting.

Across all GRPO configurations, reward and rel\_F1 show strong negative correlation ($r{=}{-}0.68$), while reward and evi\_F1 are strongly positively correlated ($r{=}0.82$). The decoupling onset occurs at step ${\sim}$30 for G=8 versus step ${\sim}$60 for G=4, consistent with faster optimization under lower-variance advantage estimates.

\subsection{Data Curation vs.\ Preference/Reward Learning}

\begin{table}[t]
\centering
\small
\begin{tabular}{llcccc}
\toprule
\textbf{Method} & \textbf{Type} & \textbf{rel\_F1} & \textbf{evi\_F1} & \textbf{EDCR} & \textbf{$\Delta$SFT} \\
\midrule
GRPO G=8 & Online reward & 7.8 & 83.5 & 92.8 & $-$32.8 \\
GRPO G=4 & Online reward & 21.7 & 57.5 & 85.4 & $-$18.9 \\
DPO+CED & Offline pref. & 34.0 & 79.8 & 66.0 & $-$6.6 \\
SFT & Supervised & 40.6 & 79.8 & 71.2 & --- \\
RSFT+flat-NLI & Data curation & 46.1 & 80.8 & 65.9 & +5.5 \\
RSFT+CED & Data curation & 46.8 & 81.6 & 63.9 & +6.3 \\
\bottomrule
\end{tabular}
\caption{Training paradigm comparison (dev set). Data curation (RSFT) succeeds where both online (GRPO) and offline (DPO) contrastive losses fail. $\Delta$SFT: rel\_F1 change in pp.}
\label{tab:spectrum}
\end{table}

Table~\ref{tab:spectrum} reveals a clear ordering across training paradigms.
DPO ($\beta{=}0.1$, 5{,}436 preference pairs) avoids GRPO's catastrophic collapse but still degrades rel\_F1 by 6.6pp below SFT.
The DPO contrastive loss simultaneously pushes up preferred and pushes down rejected candidates; in structured IE, ``rejected'' candidates may contain correct relation patterns with merely weaker evidence, so pushing them down damages extraction capability.
DPO is also more sensitive to CED label noise (${\sim}$18.5\% error rate) than RSFT, which only uses top positive candidates and never explicitly penalizes negatives.

The key insight is that the distinction is not online vs.\ offline RL, but \emph{contrastive/reward-based loss vs.\ data curation}: any method that modifies gradients based on evidence-quality comparisons risks damaging structured generation, while RSFT succeeds precisely because it is not RL---it is quality-filtered SFT.

\subsection{EDCR as Diagnostic Metric}

\begin{table}[t]
\centering
\small
\begin{tabular}{lcccc}
\toprule
\textbf{Method} & \textbf{rel\_F1} & \textbf{evi\_F1\textsubscript{c}} & \textbf{EDCR} & \textbf{Zone} \\
\midrule
GRPO G=8 & 7.8 & 83.5 & 92.8 & Collapse \\
GRPO G=4 & 21.7 & 57.5 & 85.4 & Collapse \\
SFT & 40.6 & 79.8 & 71.2 & Baseline \\
CED-reranker & 42.2 & 80.9 & 66.8 & Filtered \\
DPO+CED & 34.0 & 79.8 & 66.0 & Filtered \\
RSFT+flat-NLI & 46.1 & 80.8 & 65.9 & Filtered \\
RSFT+CED & 46.8 & 81.6 & 63.9 & Filtered \\
\bottomrule
\end{tabular}
\caption{Cross-method EDCR diagnostic (dev set). EDCR discriminates three zones: collapse ($>$85\%), baseline (${\sim}$71\%), and filtered (64--67\%).}
\label{tab:edcr_diagnostic}
\end{table}

Table~\ref{tab:edcr_diagnostic} demonstrates EDCR's ability to discriminate three method categories.
EDCR correlates negatively with rel\_F1 ($\rho{=}{-}0.52$) and evi\_F1 ($\rho{=}{-}0.38$) across methods, capturing the intuition that better methods produce fewer distractor citations.

\paragraph{Human Validation.}
We validate EDCR through a 100-sample annotation study with two annotators ($\kappa{=}0.62$).
Table~\ref{tab:edcr_human} summarizes the results.
EDCR achieves Pearson $r{=}0.44$ and Spearman $\rho{=}0.41$ against human judgments at the sample level.
While modest, method-level rank concordance is perfect (1.0), confirming EDCR reliably discriminates between methods.
EDCR systematically overestimates the distractor rate by 8.5pp, primarily because 14\% of sentences labeled as distractors by the automatic metric are judged by annotators as actually supporting the relation---reflecting incomplete gold evidence annotations in DocRED rather than metric failure.

\begin{table}[t]
\centering
\small
\begin{tabular}{lc}
\toprule
\textbf{Statistic} & \textbf{Value} \\
\midrule
Samples / Annotators & 100 / 2 \\
Inter-annotator $\kappa$ & 0.62 \\
EDCR vs.\ human Pearson $r$ & 0.44 \\
EDCR vs.\ human Spearman $\rho$ & 0.41 \\
EDCR overestimate & +8.5pp \\
False distractor rate & 14\% \\
Method-level rank concordance & 1.0 \\
\bottomrule
\end{tabular}
\caption{EDCR human validation statistics.}
\label{tab:edcr_human}
\end{table}

\subsection{CED as Inference-Time Reranker}

Without any additional training, CED can be used to rerank $N{=}8$ SFT candidates at inference time.
CED-reranker achieves rel\_F1=42.2\% (+1.6pp over SFT) and EDCR=66.8\% ($-$4.4pp over SFT) on 985 documents, but at 8.5$\times$ inference cost.
Compared to RSFT's +6.3pp rel\_F1 improvement at no additional inference cost, reranking demonstrates that CED signals are more effectively transmitted through training data curation than inference-time selection.

\subsection{Top-$k$ Selection Ablation}

\begin{table}[t]
\centering
\small
\begin{tabular}{ccccc}
\toprule
\textbf{top-$k$} & \textbf{$N_\text{train}$} & \textbf{rel\_F1} & \textbf{evi\_F1} & \textbf{EDCR} \\
\midrule
1 & 2{,}718 & \textbf{46.8} & \textbf{81.6} & \textbf{63.9} \\
3 & 8{,}154 & 45.8 & 80.2 & 66.2 \\
\bottomrule
\end{tabular}
\caption{Top-$k$ selection ablation (dev set). Strict filtering (top-1) outperforms relaxed filtering (top-3) despite 3$\times$ less training data.}
\label{tab:topk}
\end{table}

Table~\ref{tab:topk} shows that top-$k{=}1$ (strict filtering, 2{,}718 samples) uniformly outperforms top-$k{=}3$ (relaxed filtering, 8{,}154 samples).
The mean CED score drops from 0.668 to 0.512 when including rank-2/3 candidates, and the lower-quality candidates introduce distractor citation patterns that degrade evidence quality (EDCR +2.3pp).
This validates the RSFT design principle: quality of curated data matters more than quantity.

\subsection{Model Size Ablation}

\begin{table}[t]
\centering
\small
\begin{tabular}{lcccc}
\toprule
\textbf{Model} & \textbf{Method} & \textbf{rel\_F1} & \textbf{evi\_F1\textsubscript{c}} & \textbf{$\Delta$rel\_F1} \\
\midrule
Qwen3-1.7B & SFT & 32.0 & 77.5 & --- \\
 & RSFT+CED & 34.8 & 79.2 & +2.8pp \\
\midrule
Qwen3-4B & SFT & 40.6 & 79.8 & --- \\
 & RSFT+CED & 46.8 & 81.6 & +6.3pp \\
\midrule
Qwen3-8B & SFT & 44.8 & 81.0 & --- \\
 & RSFT+CED & 48.8 & 82.5 & +4.0pp \\
\bottomrule
\end{tabular}
\caption{Model size ablation (dev set). RSFT consistently improves across 1.7B, 4B, and 8B.}
\label{tab:model_size}
\end{table}

Table~\ref{tab:model_size} demonstrates consistent RSFT improvement across three model sizes (+2.8/+6.3/+4.0pp).
Smaller models generate fewer high-quality candidates, reducing RSFT training data and limiting gains.
The largest absolute performance is Qwen3-8B RSFT (rel\_F1=48.8\%), though a gap to LMRC (60.7\%, 8B, no evidence) remains, reflecting the additional burden of evidence generation.

\subsection{Re-DocRED Generalization}

\begin{table}[t]
\centering
\small
\begin{tabular}{lcccc}
\toprule
\textbf{Method} & \textbf{DocRED} & \textbf{Re-DocRED} & \textbf{$\Delta$SFT} & \textbf{EDCR} \\
\midrule
SFT & 40.6 & 48.5 & --- & 66.8 \\
RSFT+CED & 46.8 & 52.0 & +3.5pp & 61.0 \\
\bottomrule
\end{tabular}
\caption{Re-DocRED evaluation (dev set, DocRED-trained models). RSFT improvement transfers (+3.5pp).}
\label{tab:redocred}
\end{table}

Models trained on DocRED and evaluated on Re-DocRED (which corrects ${\sim}$30\% false negative triples) show consistent RSFT improvement (+3.5pp, Table~\ref{tab:redocred}).
The delta narrows from +6.3pp on DocRED to +3.5pp on Re-DocRED, a ceiling effect: Re-DocRED's more complete annotations strengthen the SFT baseline, compressing the improvement margin.

\subsection{Baseline Reproduction}

We reproduce DREEAM and SAIS on the DocRED dev set to validate evaluation protocol alignment (Table~\ref{tab:baseline_repro}).

\begin{table}[t]
\centering
\small
\begin{tabular}{lccc}
\toprule
\textbf{Method} & \textbf{rel\_F1} & \textbf{Evi-F1\textsubscript{j}} & \textbf{Evi-F1\textsubscript{c}} \\
\midrule
DREEAM (repro.) & 62.5 & 58.1 & ${\sim}$93 \\
SAIS (repro.) & 60.8 & 56.2 & ${\sim}$92 \\
RSFT+CED (ours) & 46.8 & 37.2 & 81.6 \\
SFT (ours) & 40.6 & 30.8 & 79.8 \\
\bottomrule
\end{tabular}
\caption{Dev set comparison with reproduced baselines (DREEAM-aligned joint Evi-F1 protocol). Reproduction gap $<$1pp from published test results.}
\label{tab:baseline_repro}
\end{table}

Reproduced dev results are within 1pp of published test results (DREEAM: 58.1 vs.\ 57.3; SAIS: 56.2 vs.\ 55.7), confirming evaluation alignment.
The joint Evi-F1 gap (37.2 vs.\ 58.1) decomposes as: joint $\approx$ recall $\times$ conditional $\approx 0.47 \times 0.82 = 0.385$, showing the gap is driven by extraction recall rather than per-triple evidence quality.

\subsection{NLI Model Ablation}

\begin{table}[t]
\centering
\small
\begin{tabular}{lcccc}
\toprule
\textbf{NLI Scorer} & \textbf{Disc.\ Acc.} & \textbf{rel\_F1} & \textbf{evi\_F1} & \textbf{EDCR} \\
\midrule
DeBERTa-v3-base & 81.5 & 46.8 & 81.6 & 63.9 \\
DeBERTa-v3-large & 84.8 & 47.2 & 82.0 & 63.2 \\
\bottomrule
\end{tabular}
\caption{NLI model ablation. Upgrading the scorer (+3.3pp discrimination accuracy) yields minimal downstream gains (+0.4pp rel\_F1), indicating the pipeline bottleneck is candidate generation quality, not scorer precision.}
\label{tab:nli_ablation}
\end{table}

Upgrading the NLI scorer from DeBERTa-v3-base to DeBERTa-v3-large improves CED discrimination accuracy by 3.3pp (81.5\%$\to$84.8\%) but yields only +0.4pp downstream rel\_F1 (Table~\ref{tab:nli_ablation}).
The pipeline bottleneck is candidate generation quality (fixed at $N{=}8$ SFT samples), not scorer precision; base-level accuracy (81.5\%) already exceeds the threshold for effective filtering.

\section{Discussion}

\paragraph{Why RSFT Succeeds Where RL Fails.}
The consistent failure of reward-based methods (GRPO, DPO) and success of data curation (RSFT) reveals a structural property of evidence-generating IE: the evidence-quality objective is misaligned with the extraction objective at the gradient level.
Reward-based methods transmit this misalignment directly through policy gradients, while RSFT insulates the model by converting the quality signal into a data selection criterion.
This connects to the theoretical result that RSFT is equivalent to PPO under certain conditions \cite{dong2024raft}---but with a crucial practical advantage: RSFT eliminates the risk of reward overoptimization because the model never optimizes a reward function.

\paragraph{The Role of CED.}
Our ablation shows CED's contrastive margin provides consistent but statistically non-significant improvements over flat-NLI scoring ($p{=}0.12$ for evi\_F1).
This is attributable to three factors: (1) DeBERTa-v3-base's positional bias causes contrastive and non-contrastive scores to converge for long documents; (2) incomplete DocRED gold evidence (${\sim}$14\% false distractors) dilutes the hard-negative signal; (3) the indirect RSFT pipeline (scorer $\to$ filter $\to$ training data $\to$ fine-tune) attenuates scoring-level differences.
Importantly, CED's top-1 vs.\ top-3 selection advantage (Table~\ref{tab:topk}) demonstrates that the \emph{ranking} quality of the scorer, even if not differentiated from flat-NLI, enables effective data curation.

\paragraph{EDCR Limitations.}
While EDCR achieves perfect method-level rank concordance with human judgments, its sample-level correlation ($r{=}0.44$) is moderate.
The 14\% false distractor rate---where automatically labeled distractors are judged by humans as valid evidence---reflects DocRED's annotation incompleteness rather than metric design failure.
EDCR should be used as a complementary diagnostic alongside Evi-F1, not as a standalone metric.

\section{Limitations}

\textbf{Absolute performance gap.} Our best model (Qwen3-8B RSFT, rel\_F1=48.8\%) trails encoder methods (DREEAM 62.8\%) and LLM methods without evidence (LMRC 60.7\%).
Evidence generation imposes an additional burden on the model; closing this gap while maintaining evidence quality is future work.

\textbf{Single NLI backbone.} CED relies on DeBERTa NLI models, which have known biases including lexical overlap sensitivity and positional bias on long documents.
Alternative evidence scorers (e.g., cross-encoder models, LLM-based evaluation) may better differentiate contrastive from non-contrastive scoring.

\textbf{GRPO scope.} Our decoupling findings are established for GRPO+QLoRA on document-level IE.
We do not claim this generalizes to full-rank GRPO, other online RL algorithms (PPO, REINFORCE), or other structured generation tasks.

\textbf{Gold annotation quality.} DocRED gold evidence annotations are incomplete (${\sim}$14\% false distractors by our human study), which affects both training (RSFT candidates scored against noisy gold) and evaluation (EDCR systematically overestimates by 8.5pp).

\textbf{Dev-to-test degradation.} RSFT shows 3.6pp dev-to-test degradation in rel\_F1 (vs.\ 0.6pp for SFT), suggesting indirect overfitting through the CED scoring pathway.
The improvement remains significant on test ($p{<}0.01$) but this asymmetric degradation warrants investigation.

\section{Conclusion}

We introduced \method{}, the first LLM-based method for document-level relation extraction that jointly generates relation triples and verifiable sentence-level evidence.
Through systematic comparison of training paradigms, we established that evidence-quality signals are effective as offline data curation criteria (RSFT, +3.2pp test rel\_F1) but destructive as online rewards (GRPO) or offline contrastive losses (DPO).
We characterized evidence-extraction decoupling in GRPO+QLoRA, showing it accelerates with larger group sizes, and introduced EDCR as a human-validated diagnostic for evidence quality.
Our work opens the path toward trustworthy, evidence-grounded document-level IE with large language models.

\begin{thebibliography}{20}

\bibitem[\protect\citeauthoryear{Yao et al.}{2019}]{yao2019docred}
Yuan Yao, Deming Ye, Peng Li, Xu Han, Yankai Lin, Zhenghao Liu, Zhiyuan Liu, Lixin Huang, Jie Zhou, and Maosong Sun.
\newblock DocRED: A large-scale document-level relation extraction dataset.
\newblock In \emph{Proc. ACL}, 2019.

\bibitem[\protect\citeauthoryear{Zhou et al.}{2021}]{zhou2021atlop}
Wenxuan Zhou, Kevin Huang, Tengyu Ma, and Jing Huang.
\newblock Document-level relation extraction with adaptive thresholding and localized context pooling.
\newblock In \emph{Proc. AAAI}, 2021.

\bibitem[\protect\citeauthoryear{Ma et al.}{2023}]{ma2023dreeam}
Youmi Ma, An Wang, and Naoaki Okazaki.
\newblock DREEAM: Guiding attention with evidence for improving document-level relation extraction.
\newblock In \emph{Proc. EACL}, 2023.

\bibitem[\protect\citeauthoryear{Xiao et al.}{2022}]{xiao2022sais}
Chaojun Xiao, Yuan Yao, Ruobing Xie, Xu Han, Zhiyuan Liu, Maosong Sun, Fen Lin, and Leyu Lin.
\newblock SAIS: Supervising and augmenting intermediate steps for document-level relation extraction.
\newblock In \emph{Proc. NAACL}, 2022.

\bibitem[\protect\citeauthoryear{Xu et al.}{2024}]{xu2024lmrc}
Wang Xu, Kehai Chen, and Tiejun Zhao.
\newblock Document-level relation extraction with large language models.
\newblock In \emph{Proc. EMNLP}, 2024.

\bibitem[\protect\citeauthoryear{Gao et al.}{2023}]{gao2023scaling}
Leo Gao, John Schulman, and Jacob Hilton.
\newblock Scaling laws for reward model overoptimization.
\newblock In \emph{Proc. ICML}, 2023.

\bibitem[\protect\citeauthoryear{Song et al.}{2025}]{song2025longrlvr}
Yige Song, Xinyuan Lu, Zhiling Luo, and Dongkuan Xu.
\newblock LongRLVR: Long-context reasoning with reinforcement learning and verifiable rewards.
\newblock In \emph{Proc. ICLR}, 2026.

\bibitem[\protect\citeauthoryear{Dong et al.}{2024}]{dong2024raft}
Hanze Dong, Wei Xiong, Deepanshu Goyal, Yihan Zhang, Winnie Chow, Rui Pan, Shizhe Diao, Jipeng Zhang, Kashun Shum, and Tong Zhang.
\newblock RAFT: Reward ranked finetuning for generative foundation model alignment.
\newblock \emph{Transactions on Machine Learning Research}, 2024.

\bibitem[\protect\citeauthoryear{Tan et al.}{2022}]{tan2022redocred}
Qingyu Tan, Lu Xu, Lidong Bing, Hwee Tou Ng, and Sharifah Mahani Aljunied.
\newblock Revisiting DocRED -- addressing the false negative problem in relation extraction.
\newblock In \emph{Proc. EMNLP}, 2022.

\end{thebibliography}

\end{document}
